% Section 1.8, from lectures/L01/appendix/b_exercises.md.

\section{Exercises}
\label{c1:sec:exercises}

\begin{aside}
\textbf{How to check your work.} Exercise~\ref{c1:ex:sop} is a short pen-and-paper warm-up on the
notation the book uses throughout; it needs neither VHDL, GHDL nor a board, and it states what a
good answer covers.

From Exercise~\ref{c1:ex:gates} onwards, every exercise that asks you to write a VHDL module comes
with a self-checking testbench under \file{lectures/L01/exercises/}. Write your module in its
directory \file{exercises/<module>/}, with the entity name and the \textbf{port order} the exercise
states, and then run it with GHDL \textemdash{} see \secref{c2:sec:testbenches} for the three
commands.

No FPGA board is needed for any exercise. The Quartus synthesis and board programming steps are
demonstrated during the session; your job afterwards is to get the VHDL right, and the testbench is
how you confirm it.
\end{aside}

% ------------------------------------------------------------------------------------------
\exerciseset{Truth tables and gates}

\exercise{A function that depends on less than it looks like it does}{Circuit}
\label{c1:ex:sop}
A function of two inputs \code{A}, \code{B} and one output \code{X} has the truth table below:

\begin{narrowtable}{ccc}
  \thead{A} & \thead{B} & \thead{X} \\
  \midrule
  0 & 0 & 1 \\
  0 & 1 & 0 \\
  1 & 0 & 1 \\
  1 & 1 & 0 \\
\end{narrowtable}

\xpart{a} Derive the sum-of-products equation for \code{X} straight out of the truth table
(\secref{c1:sec:boolean}). Two rows have $X = 1$, so you get two AND terms.

\xpart{b} Simplify it algebraically. Factor out what the two terms share, and apply $A + A' = 1$
and $1 \cdot B' = B'$ from \secref{c1:sec:boolean}. You should end up with a single literal.

\xpart{c} Your answer to \textbf{b)} says something about \code{A}. What? Which single gate from
the table in \secref{c1:sec:gates} computes the function, and what happens to the input \code{A}
when you draw it?

\xpart{d} Build that single-gate version in CircuitVerse and confirm that it reproduces every row
of the truth table above. Leave the input \code{A} unconnected, or simply leave it out of the
drawing; \textbf{b)} is the justification for doing so.

% ------------------------------------------------------------------------------------------
\exerciseset{Entities and architectures}

\exercise{\code{and_gate} and \code{nand_gate}}{VHDL}
\label{c1:ex:gates}
Write two small gates as complete VHDL modules, \code{and_gate} and \code{nand_gate}.

Neither gate is the point, and you should not expect either of them to be hard. The point is that
this is where an entity and an architecture stop being something you read about in
\secref{c1:sec:vhdl} and become something you write, and that the two halves are worth writing in
that order.

Both entities have the same ports:

\begin{narrowtable}{lll}
  \thead{Port} & \thead{Direction} & \thead{Type} \\
  \midrule
  \code{a} & in & \code{std_logic} \\
  \code{b} & in & \code{std_logic} \\
  \code{x} & out & \code{std_logic} \\
\end{narrowtable}

\xpart{a} Write the \code{entity} declaration for \code{and_gate}, and nothing else. In VHDL a
module's interface is something you can write down and reason about before any implementation
exists at all, which is the habit this part exists to build. Identify which ports are required, the
direction of each, and the type of each.

\xpart{b} Add an \code{architecture} that implements the AND gate with a single concurrent signal
assignment.

\xpart{c} Now write \code{nand_gate}: the same interface with an inverted output, using VHDL's
\code{nand} operator directly. Do \textbf{not} write \code{not (a and b)}.

The two are equivalent as Boolean algebra, and on an FPGA they make no difference at all, since
both end up in the same lookup table. It matters on the kind of hardware Exercise~\ref{c1:ex:nand}
below is about, where NAND is the gate you physically have and everything else must be built from
it. Writing what you mean, and letting the tools decide what it costs, is the habit; reaching past
an operator the language already gives you is what to avoid.

\modulefig[0.52]{lectures/L01/appendix/images/and_gate.png}
\modulefig[0.52]{lectures/L01/appendix/images/nand_gate.png}

\begin{selfcheck}
\textbf{Self-check.} Name your entities \code{and_gate} and \code{nand_gate}, each with the ports
\code{a}, \code{b} (in) and \code{x} (out), declared in that order; their testbenches are in
\file{exercises/and_gate/} and \file{exercises/nand_gate/}, and each checks all four rows of the
corresponding truth table.
\end{selfcheck}

\exercise{\code{xor3}: odd parity with three inputs}{VHDL}
\label{c1:ex:xor3}
An XOR gate has three inputs, \code{a}, \code{b}, \code{c}, and one output, \code{x}. The output is
\code{1} when an odd number of the inputs are \code{1} (three-input XOR, or the parity function).

\xpart{a} Write the truth table for \code{x} over all eight combinations of \code{a}, \code{b},
\code{c}, and confirm that it matches the description ``an odd number of \code{1}s''.

\xpart{b} Implement the gate in VHDL, as a module of its own with an entity whose ports are three
\code{std_logic} inputs and one \code{std_logic} output, and an architecture that drives \code{x}
with a single concurrent assignment. \Secref{c1:sec:vhdl} shows the entity/architecture pattern and
notes that VHDL's \code{xor} operator works directly on \code{std_logic}; \secref{c1:sec:module}
shows a complete module from start to finish.

\xpart{c} Build the same gate in CircuitVerse and confirm that it reproduces every row of your
truth table.

\note[Hint]\code{xor} chains from left to right, so the whole architecture body is a single line,
\code{x <= a xor b xor c;}, with no need for either an internal \code{signal} or intermediate
gates.

\modulefig[0.52]{lectures/L01/appendix/images/xor3.png}

\begin{selfcheck}
\textbf{Self-check.} Name your VHDL entity \code{xor3}, with the ports \code{a}, \code{b},
\code{c} (in) and \code{x} (out), declared in that order; its testbench is in
\file{exercises/xor3/}, and \secref{c1:sec:module} lists the three commands that run it.
\end{selfcheck}

\exercise{\code{majority3}: a triple-redundant vote}{VHDL}
\label{c1:ex:majority3}
A \textbf{majority gate} has three inputs, \code{a}, \code{b}, \code{c}, and one output, \code{x}.
The output is \code{1} when \emph{at least two} of the three inputs are \code{1}.

This is the circuit behind a triple-redundant vote: three sensors report the same measurement, and
the system acts on the answer two of them agree on, so that a single broken sensor cannot decide
anything on its own.

\xpart{a} Write the truth table for \code{x} over all eight combinations of \code{a}, \code{b},
\code{c}.

\xpart{b} Derive the sum-of-products equation for \code{x} straight out of the truth table
(\secref{c1:sec:boolean}):
\begin{itemize}
  \item Four rows have $x = 1$, so you get four AND terms.
  \item Every term contains all three inputs, primed where the input is \code{0} on its row.
\end{itemize}

\xpart{c} Implement the gate in VHDL, as a module of its own:
\begin{itemize}
  \item An entity with three \code{std_logic} inputs and one \code{std_logic} output.
  \item An architecture that drives \code{x} with a single concurrent assignment, translating your
    four AND terms directly (\secref{c1:sec:vhdl}).
\end{itemize}

\xpart{d} Build the same network in CircuitVerse and confirm that it reproduces every row of your
truth table.

\note[Hint]Write the equation out on paper before you open either the editor or CircuitVerse. Four
three-input AND terms plus a four-input OR gate is a lot of wiring to correct after the fact.

\note[Looking ahead]Count the gates your equation calls for, and then convince yourself by
inspection that $x = ab + ac + bc$ computes the same function with three two-input AND gates
instead of four three-input ones. Sum of products gave you a \emph{correct} network, not the
\emph{smallest}; finding the smaller one systematically, rather than by staring at it, is what
\chapref{c2:ch} is for.

\modulefig[0.52]{lectures/L01/appendix/images/majority3.png}

\begin{selfcheck}
\textbf{Self-check.} Name your VHDL entity \code{majority3}, with the ports \code{a}, \code{b},
\code{c} (in) and \code{x} (out), declared in that order; its testbench is in
\file{exercises/majority3/} and checks all eight rows. Either equation passes \textemdash{} the
testbench checks the function, not which network you chose.
\end{selfcheck}

\exercise{\code{half_adder}: the first module with two outputs}{VHDL}
\label{c1:ex:half_adder}
Write a \textbf{half adder}: \code{sum} and \code{carry} from two single bits. You know the
circuit; the reason it is here is that it is the book's first module with \textbf{two outputs}, and
that is a VHDL question rather than a logic question.

\xpart{a} Implement it, with two concurrent assignments and no internal \code{signal}.

Two outputs do not mean two entities, two architectures or a process. Every output port is driven
by its own concurrent assignment, and the two are not steps that happen in an order: they describe
two separate pieces of wire, both live all the time. If you catch yourself reaching for a process
in order to ``return'' two values, that instinct is the software instinct, and
\secref{c1:sec:vhdl}'s ``a concurrent assignment describes a wire'' is the correction.

\xpart{b} Build it in CircuitVerse and confirm both outputs, so that you have seen the two-output
shape as a drawing before you meet it as an entity with several \code{out} ports in
\chapref{c2:ch}.

\note[Looking ahead]This is the cell the counter in \chapref{c6:ch} is built from. The four-bit
adder there is four of these chained carry-to-carry, and the counter is that adder with its output
fed back to its input, which is why the wraparound needs no logic at all.

\modulefig[0.56]{lectures/L01/appendix/images/half_adder.png}

\begin{selfcheck}
\textbf{Self-check.} Name your entity \code{half_adder}, with the ports \code{a}, \code{b} (in) and
\code{sum}, \code{carry} (out), declared in that order; its testbench is in
\file{exercises/half_adder/} and checks \code{sum} and \code{carry} separately, so that a design
that has one right and the other wrong is told which is which.
\end{selfcheck}

% ------------------------------------------------------------------------------------------
\exerciseset{Universal gates}

\exercise{\code{or_from_nand}}{VHDL}
\label{c1:ex:nand}
\Secref{c1:sec:gates} states that every Boolean function can be built from \code{NAND} gates
alone. This exercise has you prove it to yourself for the three operators you have been using.

\xpart{a} Derive the \code{OR} case on paper before you build anything, since De Morgan
(\secref{c1:sec:boolean}) gives it to you directly. Start from $(A'B')'$ and apply the law; you
should arrive at $A + B$. Now read that expression back as a circuit: what is $A'B'$ with an
inversion on its output, and how many \code{NAND} gates is the whole thing?

\xpart{b} Build each of the following in CircuitVerse from \code{NAND} gates alone, checking each
against its truth table in \secref{c1:sec:gates} as you go:
\begin{itemize}
  \item \code{NOT a}: one NAND gate is enough, if you feed the same signal to both its inputs.
  \item \code{a AND b}: one NAND gate, followed by the \code{NOT} you just built.
  \item \code{a OR b}: the network you derived in \textbf{a)}. Check your gate count against what
    you predicted.
\end{itemize}

\xpart{c} Write the \code{OR} version in VHDL, as a module of its own:
\begin{itemize}
  \item An entity with two \code{std_logic} inputs and one \code{std_logic} output.
  \item An architecture that drives \code{x} with a single concurrent assignment.
  \item Use VHDL's \code{nand} operator alone: no \code{or}, no \code{and}, no \code{not}
    anywhere in the expression.
\end{itemize}

\xpart{d} Explain why this matters in practice. Bear in mind that a chip manufacturing process has
to physically implement every gate type it offers, and that \code{NAND} is the cheapest gate to
build in CMOS.

\note[Hint]VHDL's \code{nand} operator does not chain the way \code{and} and \code{or} do:
\code{a nand b nand c} is not valid. Parenthesise every NAND explicitly, which is what you want
here anyway, since every pair of parentheses is a physical gate in your CircuitVerse drawing.

\modulefig[0.52]{lectures/L01/appendix/images/or_from_nand.png}

\begin{selfcheck}
\textbf{Self-check.} Name your VHDL entity \code{or_from_nand}, with the ports \code{a}, \code{b}
(in) and \code{x} (out), declared in that order; its testbench is in
\file{exercises/or_from_nand/} and checks it against all four rows of the \code{OR} truth table.
\end{selfcheck}

% ------------------------------------------------------------------------------------------
\exerciseset{Putting it all together}

\exercise{The brake assistant in VHDL}{VHDL}
\label{c1:ex:adas}
Write the brake assistant from \secref{c1:sec:adas} in VHDL.

\Secref{c1:sec:adas} had you derive its equations, count its gates and sketch its network
\emph{before} the session, and the session then built it live. This is where you write it yourself,
and where a testbench decides whether your reading of the requirement matches the one in the truth
table.

The entity must have:

\begin{booktable}{@{}p{8em}p{3.4em}p{5.6em}L@{}}
  \thead{Port} & \thead{Dir.} & \thead{Type} & \thead{Description} \\
  \midrule
  \code{driver_brake} & in & \code{std_logic} & The driver is on the brake pedal. \\
  \code{sensor} & in & \code{std_logic} & The distance sensor sees an obstacle. \\
  \code{radar} & in & \code{std_logic} & The radar sees an obstacle. \\
  \code{error} & in & \code{std_logic} & The assistance system reports a fault. \\
  \code{engine_brake} & out & \code{std_logic} & Apply the brakes. \\
\end{booktable}

\xpart{a} Implement it from \textbf{your own} equations rather than from what was on the screen. If
the two disagree, that disagreement is the most useful thing this exercise can give you, and the
testbench says which of them the truth table sides with.

Use two concurrent assignments and one internal \code{signal}:
\begin{itemize}
  \item The signal holds the assistance system's own decision, before the driver's pedal is
    weighed in.
  \item This is the book's first exercise with an internal signal, and it is the reason the keyword
    exists: a name for an intermediate result, so that every line still reads as the sentence in
    the requirement it came from.
\end{itemize}

\xpart{b} Now break it deliberately, in the specific way \secref{c1:sec:adas} warns about. Feed
\code{driver_brake} \emph{through} the fault logic instead of letting it go past:

\begin{vhdlcode}
engine_brake <= (driver_brake or sensor or radar) and (not error);
\end{vhdlcode}

It is one gate simpler, and it looks reasonable. Run the testbench against it and answer:
\begin{itemize}
  \item Which input combination does it fail on first, and what happens physically in the car at
    that moment?
  \item How many of the sixteen rows does it get wrong, and what do they all have in common?
  \item The rows that fail are exactly the ones the safety requirement was written for. What does
    that tell you about the value of a truth table derived from a requirement, as opposed to one
    read off a circuit you have already built?
\end{itemize}

\modulefig[0.72]{lectures/L01/appendix/images/adas.png}

\begin{selfcheck}
\textbf{Self-check.} Name your entity \code{adas}, with the inputs \code{driver_brake},
\code{sensor}, \code{radar}, \code{error} and the output \code{engine_brake}, declared in that
order; its testbench is in \file{exercises/adas/} and checks all sixteen rows of
\secref{c1:sec:adas}'s truth table.
\end{selfcheck}
